\documentclass{article}
\usepackage[utf8]{inputenc}
\usepackage[russian]{babel}
\usepackage{amsmath}
\usepackage{amssymb}
\usepackage{amsthm}
\usepackage{hyperref}

\title{Краткое изложение пяти статей по термоупругости}
\date{\today}

\begin{document}
\maketitle
\tableofcontents

\section*{Краткое изложение статей}

\subsection*{1. The Propagation of Thermoelastic Waves in Different Anisotropic Media Using Matricant Method}

Данная работа фокусируется на анализе \textbf{распространения продольных упругих и тепловых волн в анизотропных средах} (моноклинной, тригональной, гексагональной и кубической кристаллических систем) [1-3].
\begin{itemize}
    \item \textbf{Методология:} Использовался \textbf{аналитический метод матриканта} [1, 4-6], который позволяет получать точные аналитические решения систем дифференциальных уравнений, описывающих связанные термоупругие процессы [5]. Этот метод обеспечивает унифицированное описание волнового распространения в различных средах [7].
    \item \textbf{Проблема:} Распространение термоупругих волн в анизотропных средах изучено не полностью, в отличие от изотропных сред [1].
    \item \textbf{Результаты:} Получено решение задачи распространения тепловых волн в одномерном случае, которое совпадает с известным классическим решением [1, 3, 8]. Определены \textbf{коэффициент затухания ($q$)} и \textbf{фазовая скорость ($\nu$)} тепловых волн для различных материалов (например, кварца, алюминия, меди) [9, 10]. Показано, что $q$ и $\nu$ зависят от угловой частоты ($\omega$) одинаковым образом, возрастая по параболическому закону ($\sim \sqrt{\omega}$) [11].
\end{itemize}

\subsection*{2. Reflection and transmission of thermoelastic waves in multilayered media}

В этом исследовании разработан обобщенный алгоритм \textbf{матрицы переноса (TM)} для расчета коэффициентов отражения/прохождения ($R/T$) термоупругих волн в многослойных средах с плоскими границами раздела [12].
\begin{itemize}
    \item \textbf{Модель:} Используется \textbf{обобщенная теория термоупругости Лорда-Шульмана (Lord-Shulman, LS)} [12-14]. Теория LS вводит время релаксации ($\tau$), что делает уравнение теплопроводности гиперболическим, предсказывая конечную скорость распространения тепла и избегая проблем классической теории Фурье [14-16].
    \item \textbf{Волны:} Модель предсказывает три типа волн: \textbf{быстрая продольная волна ($P$), сдвиговая волна ($S$) и тепловая волна ($T$)} (или медленная P-волна), которая возникает благодаря связи между деформацией и температурой [15, 17, 18].
    \item \textbf{Верификация:} Коэффициенты $R/T$ проверяются с помощью \textbf{сохранения энергии} (energy balance/conservation of energy) [12, 17, 19].
    \item \textbf{Применение:} Модель применяется к геотермальной системе \textbf{горячих сухих пород (Hot Dry Rock, HDR)} [12]. Более высокая абсолютная температура ($T_0$) приводит к увеличению амплитуд отраженных и проходящих $S$-волн [20, 21]. Коэффициенты $R/T$ могут быть использованы для получения информации о тепловых характеристиках многослойной среды [12, 21].
\end{itemize}

\subsection*{3. Distribution of coupled thermoelastic waves in the different rocks using MatLab GUI application}

Работа описывает создание и использование компьютерного приложения \textbf{GUI MATLAB} для численного и графического анализа волновых чисел и скоростей в связанных термоупругих системах [22-24].
\begin{itemize}
    \item \textbf{Цель:} Изучить влияние частоты ($\omega$) на поведение первичных ($P$) и вторичных ($S$) волн, используя физические характеристики горных пород из месторождений Казахстана (гранит, песчаник, алевролит, сланец, известняк) [22, 25-27].
    \item \textbf{Математическая модель:} Динамика среды описывается системой дифференциальных уравнений \textbf{смешанного гиперболическо-параболического типа} [28, 29]. Для решения использовался метод граничных интегральных уравнений с обобщенными функциями [30, 31].
    \item \textbf{Волновые числа:} Характеристическое уравнение имеет \textbf{шесть корней (волновых чисел)} [32]:
    \begin{itemize}
        \item Четыре комплексных корня ($\zeta_1^2, \zeta_2^2$) связаны с \textbf{$P$-волнами} (продольными, волнами сжатия-растяжения), которые распространяются параллельно направлению движения волны и могут проходить через твердые тела, жидкости и газы [33-35].
        \item Два вещественных корня ($\zeta_3^2$) связаны с \textbf{$S$-волнами} (поперечными, сдвиговыми), которые колеблются перпендикулярно направлению распространения и могут распространяться \textbf{только в твердых материалах} [33, 36, 37].
    \end{itemize}
    \item \textbf{Значимость:} Приложение предоставляет информацию о влиянии геологических условий на распространение волн, что важно для оценки сейсмического риска и проектирования устойчивых сооружений [22, 27].
\end{itemize}

\subsection*{4. Amplitude-variation-with-offset in thermoelastic media}

В статье разработана аппроксимация \textbf{изменения амплитуды с удалением (AVO)}, основанная на теории термоупругости \textbf{Лорда-Шульмана (LS)} [38].
\begin{itemize}
    \item \textbf{Модель:} Термоупругая модель предсказывает две волны сжатия ($P$ и $T$) и одну сдвиговую волну ($S$) [38]. Аппроксимация AVO включает термические эффекты и основана на линеаризованном приближении Аки-Ричардса (Aki-Richards) [39, 40].
    \item \textbf{Верификация:} Точные коэффициенты отражения/прохождения ($R/T$) рассчитываются с использованием обобщенного закона Снелла и подтверждаются \textbf{сохранением энергии} [39, 41, 42]. Аппроксимация AVO была проверена сравнением с синтетическими сейсмограммами и реальными данными месторождения Таримского бассейна, где температура превышает $410 \, \text{K}$ [38, 43, 44].
    \item \textbf{Влияние температуры ($T_0$):} Увеличение абсолютной температуры ($T_0$) в среде прохождения приводит к \textbf{уменьшению амплитуды отраженной $P$-волны} ($R_{PP}$) и \textbf{увеличению амплитуды проходящей $P$-волны} ($T_{PP}$) из-за снижения контраста импеданса [43].
    \item \textbf{Точность:} Термоупругая аппроксимация AVO дает \textbf{более точное время прохождения} (traveltimes) на малых углах по сравнению с традиционным упругим AVO [45].
\end{itemize}

\subsection*{5. Analytical wave solutions in thermoelastic media with temperature-dependent properties via IMETF method}

В этой работе представлено аналитическое исследование точных волновых решений, полученных в контексте обобщенной модели термоупругости с \textbf{тремя фазовыми запаздываниями (Three-Phase-Lag, 3PHL)} [46, 47].
\begin{itemize}
    \item \textbf{Нелинейность и модель:} Ключевой особенностью является явное включение \textbf{зависимости свойств материала от температуры}, что вносит нелинейность в управляющие уравнения [46, 48, 49]. Модель 3PHL вводит три задержки ($\tau_q, \tau_\theta, \tau_\nu$) для учета нефурьевой теплопроводности и связи между тепловым и механическим откликами [47].
    \item \textbf{Методология:} Используется метод \textbf{Improved Modified Extended Tanh Function (IMETF)} — мощный аналитический инструмент для получения точных решений нелинейных дифференциальных уравнений [46, 50].
    \item \textbf{Решения:} Метод IMETF позволил вывести разнообразные точные аналитические решения, включая \textbf{гиперболические, экспоненциальные, эллиптические Вейерштрасса и колоколообразные уединенные волновые решения} (bell-shaped solitary wave solutions) [46, 51].
    \item \textbf{Наблюдения (на примере меди):} Увеличение параметра чувствительности к температуре ($\alpha$) обычно приводит к \textbf{снижению интенсивности теплового поля} и \textbf{смещения}, что может указывать на то, что материал лучше рассеивает тепловую волну при более высокой температурной зависимости [52, 53]. При предельных условиях ($\alpha \to 0$ и $\tau_\nu \to 0$) решения согласуются с моделями Dual-Phase-Lag и классической термоупругости [54, 55].
\end{itemize}


\subsection*{Источники}
    \begin{itemize}
        \item \href{https://pdfs.semanticscholar.org/488a/91fa54fe48b2a361f9a4caafd9e606447f7f.pdf}{The Propagation of Thermoelastic Waves in Different Anisotropic Media Using Matricant Method}
    \end{itemize}


    \begin{itemize}
        \item \href{https://www.nature.com/articles/s41598-025-09344-w.pdf}{Analytical wave solutions in thermoelastic media with temperature-dependent properties via IMETF method}
    \end{itemize}

    
    \begin{itemize}
        \item \href{https://ricerca.ogs.it/retrieve/a16968a5-6f6f-496e-819f-d244ff797fbd/hou-et-al-2023-amplitude-variation-with-offset-in-thermoelastic-media.pdf}{Amplitude-variation-with-offset in thermoelastic media}
    \end{itemize}

    
    \begin{itemize}
        \item \href{https://ceur-ws.org/Vol-3966/W1Paper2.pdf}{Distribution of coupled thermoelastic waves in the different rocks using MatLab GUI application}
    \end{itemize}


    \begin{itemize}
        \item \href{https://academic.oup.com/gji/article/232/2/1408/6761682?login=false}{Simulation of thermoelastic wave propagation in 3-D multilayered half-space media}
    \end{itemize}
\end{document}